\summarychapter{26}{Thrust into the Game}{John Calvin}

\begin{inbrief}
John Calvin, a shy French scholar pressed into service at Geneva, shaped the third tradition of the
Reformation: Reformed or Calvinistic Christianity, to which Presbyterians, the Reformed churches of Germany and the Netherlands, and many Congregationalists and Baptists belong. The chapter traces his life, his
disciplined Geneva, his central belief in the sovereignty of God, and the spread of Calvinism to
France, the Netherlands and John Knox's Scotland.
\question{What are the distinctive features of Reformed Christianity?}
\end{inbrief}

War between France and Spain had closed the road to Strasbourg, so in July 1536 the young scholar had
to pass through Geneva, intending to stay one night. The pleasure-loving city had thrown off the Duke
of Savoy and the pope and was torn by factions. William Farel had preached there for four years and
the mass had ceased, but Geneva's Protestantism rested on political hostility to its bishop rather
than conviction, and no one had shaped its institutions on biblical lines. Farel called on the visitor
and urged him to stay. When Calvin pleaded his studies, Farel warned that God would punish him for
seeking his own interest. Terror-stricken, Calvin stayed. By nature ``a bit antisocial and shy,'' he
later said, God had thrust him into the limelight and made him ``get into the game.''

\ssection{The Meaning of Reformed Christianity}
As Luther's central doctrine was justification by faith, Calvin's was the sovereignty of God. Both
had an overwhelming sense of God's majesty; for Luther it set off the miracle of forgiveness, for
Calvin it guaranteed the certainty of God's purpose. Calvin (1509--1564) shared the four Protestant
answers but was born a generation later in another country and was a different man. Luther was a
peasant, monk and professor; Calvin a scholar and lawyer called to public ministry in a thriving
commercial city. Their foundations were the same, the structures they built at Wittenberg and Geneva
different.

Calvin built on Zwingli. The Zurich reform spread quickly through German-speaking Switzerland: Bern
in 1528, Basel in 1529, then Appenzell, St.~Gall and Schaffhausen. In Strasbourg Martin Bucer
(1491--1551), second in influence in Germany only to Luther and Melanchthon, sympathized with Zwingli.
But conflict between Protestant and Catholic cantons led on October 11, 1531, to the battle of Kappel,
in which Zwingli was killed. Heinrich Bullinger succeeded him at Zurich, and by the 1540s
French-speaking Geneva had become the international center of Reformed Christianity.

Calvin grew up in a small town of northern France, some sixty miles from Paris. He entered the University of Paris at
fourteen and acquired a brilliant style and a skill in argument, so that people might dislike what he
said but could not mistake it; he took his master's degree in 1528. At his father's wish he studied
law at Orl\'eans and Bourges, and after his father's death in 1531 he returned to Paris and the
classics. There he met reforming ideas and experienced what he called an ``unexpected conversion,''
more than enlightenment: he surrendered his will to God, gave up his scholar's career and joined the
Protestant cause. In the autumn of 1533 his friend Nicholas Cop, rector of the university, gave a
strongly Protestant address that some thought Calvin had written, and Calvin had to flee. At Basel in
March 1536 he published the first edition of the \emph{Institutes of the Christian Religion}, which no other Protestant statement of doctrine in that age matched for clarity, logic and readability. It made
him famous overnight, and he enlarged it all his life without changing its substance. Its preface, a
letter to King Francis I defending the French Protestants, made him a leader of their cause.

\ssection{Geneva: Home of Calvinism}
Appointed Professor of Sacred Scriptures, Calvin went to work at once. He drew up a confession of
faith to be accepted by every citizen, planned education for all, and insisted on excommunication,
exclusion from the Lord's Supper, for those whose lives fell short. It was the strictest moral
discipline in Protestantism and more than the city fathers had bargained for. The dispute turned on
who may excommunicate, church or magistrates. After a year Calvin and Farel lost, and in April 1538
they were ordered out. The three years that followed in Strasbourg were probably his happiest. He was
pastor of the French refugee church, taught theology, represented the city at religious conferences
in Germany and married a widow with two children, his companion until her death in March 1549.

When his friends regained power, Geneva begged him back, and in September 1541 he returned
reluctantly, almost on his own terms. The church constitution provided four offices: pastors,
teachers, elders and deacons. Twelve elders with the ministers formed the Consistory, which
supervised the city's morals, dealing with everything from absence from worship to drinking,
adultery, gambling and dancing. Opposition continued, and Calvin was often near banishment, but the
stream of religious exiles, mostly French, steadily enlarged his following. In 1553, at a low point
in his influence, the brilliant and erratic Spanish physician Michael Servetus, fleeing Catholic
persecution for denying the Trinity, arrived in Geneva. Calvin wanted a more merciful death than
burning but supported silencing him; Servetus was burned, and later generations remembered Calvin as
``the man who burned Servetus.'' Two years later his position was secure. He saw Geneva as a refuge
for the persecuted, an example of a disciplined Christian community and a training center for
ministers. Students came from all over Europe to what John Knox called ``the most perfect school of
Christ that ever was on earth since the days of the Apostles.''

\ssection{The Sovereignty of God}
What they learned flowed from God's absolute sovereignty. God is the ``Governor of all things,'' who
from eternity decreed what he would do and by his power carries it out. This is not general guidance
or fate but personal direction of individual lives: no sparrow falls unknown to the Father. If
Luther's text was ``The just shall live by faith,'' Calvin's was ``Thy will be done, on earth as it
is in heaven.'' Predestination was an old doctrine, held by Paul, Augustine and Luther, and Calvin found in it a spur to devotion: election to eternal life is the deepest ground of confidence, humility and moral
power. He did not claim to know who the elect are, but three tests were a fair measure for a
disciplined church: participation in the two sacraments, an upright life and a public profession of
faith.

More than for Luther, faith for Calvin issues in strenuous effort to bring in God's kingdom on earth.
The Christian is no longer judged by the law but finds in it the pattern of moral character; no one
is justified by works, yet no justified person is without them. This pursuit of righteousness made
character the test of real religion and explains Calvinism's social activism. Sovereignty also shaped
its view of the state. Luther tended to treat the state as supreme, and German princes decided where
and how the gospel was preached. Calvin taught that no man, pope or king, may claim absolute power. He
never preached a right of revolution, but he encouraged representative assemblies and their right to
resist tyranny, a key factor in the growth of constitutional government. The church is subject to
secular rulers only in secular matters and must guide them in spiritual ones. His followers left
Geneva, which they saw as God's beachhead, like a spiritual conspiracy against false religion and
restrictive government, and Calvinism quickly became international.

In France the Calvinists, called Huguenots, were a minority made important by converts among the
nobility. They were close to seizing leadership when thousands were massacred on St.~Bartholomew's
Day in 1572, and from then on they no longer endangered the Catholic monarchy. In the
Netherlands Calvinism rallied the resistance to Catholic Spain, with ministers among the first
leaders. William the Silent, head of the national party in the north, joined the Reformed church in
1573 and helped fashion a Dutch republic; the national anthem was written for his followers. In
Scotland the Calvinists created what no other country of the century had: a nation that held one religion under a ruler who held another.

\ssection{John Knox's Scotland}
That ruler was Mary Queen of Scots, eighteen years old, living abroad and married into the French
royal house, so that Scots and English feared she would hand Scotland to France. John Knox, a restless
activist, had worked for Calvinism in England until 1553, when the country returned to Catholicism
under Henry VIII's daughter Mary I, whose persecutions earned her the name Bloody Mary. On the
Continent he worked out the theory that Protestants may resist, by force if necessary, a Catholic
ruler who hinders their worship and mission. That went further than Calvin, but it appealed to many
Scottish nobles. When civil war broke out in 1559 Knox hurried home. By the summer of 1560 the
Calvinists held Edinburgh, and he drafted the articles of religion by which parliament abolished Roman
Catholicism. The next year Mary, now a widow of nineteen, came back to a Protestant kingdom. Knox and
the queen came to symbolize the conflict of the age: Protestant against Catholic, and the democratic
claims of Calvinism against the crown's power to appoint bishops. Events went Knox's way. Mary's heirs tried to undo his work, but no country on earth stayed more firmly Calvinist than Scotland. A reformed Geneva was the least of what Calvin left behind at his death in 1564. Across Europe, and
soon in America, his followers carried on the game.

\begin{keyfacts}
\keypeople{John Calvin; William Farel; Ulrich Zwingli; Heinrich Bullinger; Martin Bucer; Nicholas Cop;
Francis I; Michael Servetus; William the Silent; John Knox; Mary Queen of Scots; Mary I (``Bloody
Mary'').}
\keydates{%
  \kd{1509}{Birth of John Calvin}
  \kd{1531}{Zwingli killed at the battle of Kappel (October 11)}
  \kd{1536}{First edition of Calvin's \emph{Institutes}; Farel keeps him in Geneva}
  \kd{1538}{Calvin and Farel expelled from Geneva}
  \kd{1541}{Calvin returns to Geneva}
  \kd{1553}{Michael Servetus burned in Geneva}
  \kd{1560}{The Scottish parliament adopts Knox's articles of religion}
  \kd{1564}{Death of Calvin}
  \kd{1572}{Massacre of Huguenots on St.~Bartholomew's Day}}
\keyterms{Reformed (Calvinistic) Christianity; sovereignty of God; predestination and election;
\emph{Institutes of the Christian Religion}; Consistory; pastors, teachers, elders, deacons;
Huguenots; resistance to tyranny.}
\end{keyfacts}
